% Inactive mathematical fragment. Research checkpoint 469.
% Real length and area are those of the associated Riemannian metric.
% Geometric restriction and radial estimates remain explicit hypotheses.

Let $C$ be a connected compact Riemann surface. Fix a smooth Hermitian metric $h$ of area one, and use the associated real metrics for lengths and areas.

\begin{lemma}\label{lem:conformal-area-diameter}
There is $K_h>0$ such that every smooth Hermitian metric $g\leq M h$, with $M>0$ and $A=\operatorname{Area}_g(C)$, satisfies
\[
 \operatorname{diam}(C,g)
 \leq K_h\sqrt{A\bigl(1+\log(M/A)\bigr)}.
\]
Here $0<A\leq M$. In particular, for a fixed upper bound $M$, area tending to zero implies diameter tending to zero.
\end{lemma}
\begin{proof}
Replace $g$ by $g/M$. Its area is $a=A/M\leq1$, and it is bounded above by $h$. We prove the estimate for this normalized metric.

Choose a finite cover by holomorphic coordinate disks $U_i$ represented by $\{|z|<1\}$, with each coordinate defined on the larger disk $\{|z|<2\}$. Write $g=\rho_i|\mathrm{d}z|^2$ there. There is a common $L>0$ such that $0<\rho_i\leq L$ on all the closed unit disks, and
\[
 \int_{U_i}\rho_i\,\mathrm{d}\lambda_2\leq a,
\]
where $\lambda_2$ is Euclidean area. For $p,z\in U_i$, let $\ell_p(z)$ be the $g$-length of their straight coordinate segment, which stays in $U_i$. Polar coordinates about $p$ and Fubini give
\[
 \int_{U_i}\ell_p(z)\,\mathrm{d}\lambda_2(z)
 \leq 2\int_{U_i}\frac{\sqrt{\rho_i(w)}}{|w-p|}
                    \,\mathrm{d}\lambda_2(w).
\]
Indeed each ray meets the convex unit disk in an interval of length $R\leq2$. Interchanging the radial integrals multiplies the length integrand at radius $s$ by $(R^2-s^2)/2\leq2$.

For $0<a\leq1/4$, put $t=\sqrt a$. The integral on the right, without its factor two, is bounded by
\[
 \begin{split}
 \sqrt L\int_{\{|w-p|<t\}}\frac{\mathrm{d}\lambda_2(w)}{|w-p|}
 &+\left(\int_{U_i}\rho_i\,\mathrm{d}\lambda_2\right)^{1/2}
   \left(\int_{\{t<|w-p|<2\}}\frac{\mathrm{d}\lambda_2(w)}{|w-p|^2}\right)^{1/2}\\
 &\leq 2\pi\sqrt L\,t+\sqrt{2\pi a\log(2/t)}
 \leq K\sqrt{a\bigl(1+\log(1/a)\bigr)}.
 \end{split}
\]
Thus for any $p,q\in U_i$, averaging $\ell_p(z)+\ell_q(z)$ over $z\in U_i$ supplies a two-segment path of length at most $K\sqrt{a(1+\log(1/a))}$. The constant is common to the finite cover.

The overlap graph of this cover is connected because $C$ is connected. Choose points in consecutive overlaps along a simple path in this finite graph. The local estimates join any two points of $C$ through at most the number of covering disks, proving the global bound for $a\leq1/4$. For $1/4\leq a\leq1$, the comparison $g\leq h$ gives $\operatorname{diam}(C,g)\leq\operatorname{diam}(C,h)$, and increasing $K_h$ gives the same estimate. Restoring the factor $M$ scales lengths by $\sqrt M$ and proves the assertion.
\end{proof}

\begin{example}\label{ex:conformal-log-diameter-sharpness}
On every such $(C,h)$ there are smooth Hermitian metrics $g_\varepsilon\leq h$, for $\varepsilon\downarrow0$, such that
\[
 A_\varepsilon:=\operatorname{Area}_{g_\varepsilon}(C)
       \asymp\varepsilon^2\log(1/\varepsilon),
 \qquad
 \operatorname{diam}(C,g_\varepsilon)
       \asymp\varepsilon\log(1/\varepsilon).
\]
Consequently the preceding diameter bound has sharp order under its uniform upper bound. Without that bound, area can tend to zero while diameter stays bounded below by a positive constant.
\end{example}
\begin{proof}
Choose a coordinate disk $\{|z|<R\}$ with $0<R<1$, on whose closure $h=\lambda(z)|\mathrm{d}z|^2$ with $\lambda>0$ smooth. Choose a smooth cutoff $0\leq\chi\leq1$, supported in this disk and equal to one for $|z|\leq R/2$, and a constant $c>0$ with $c\leq\min\lambda$. For sufficiently small $\varepsilon>0$, set
\[
 g_\varepsilon=
 \left(\chi\frac{c\varepsilon^2}{\lambda(z)(\varepsilon^2+|z|^2)}
             +(1-\chi)\varepsilon^2\right)h
\]
in the disk, and set $g_\varepsilon=\varepsilon^2h$ outside it. The cutoff makes this a smooth positive metric, and both coefficients in its convex combination are at most one.

On $|z|\leq R/2$ its metric density is $c\varepsilon^2/(\varepsilon^2+|z|^2)$, whose area is
\[
 \pi c\varepsilon^2
       \log\left(1+\frac{R^2}{4\varepsilon^2}\right).
\]
The remaining area is $O(\varepsilon^2)$, proving the asserted area comparison. Any path from the origin to a point outside this smaller disk must first cross its boundary. Radial variation is bounded by Euclidean speed, so that initial portion has length at least
\[
 \sqrt c\,\varepsilon\int_0^{R/2}
                 \frac{\mathrm{d}r}{\sqrt{\varepsilon^2+r^2}}
       =\sqrt c\,\varepsilon\operatorname{arsinh}(R/(2\varepsilon))
       \asymp\varepsilon\log(1/\varepsilon).
\]
This gives the diameter lower bound. The upper bound follows from \cref{lem:conformal-area-diameter}, since $A_\varepsilon(1+\log(1/A_\varepsilon))\asymp\varepsilon^2\log^2(1/\varepsilon)$. Finally rescale by
\[
 \widetilde g_\varepsilon
       =\frac{g_\varepsilon}{\varepsilon^2\log^2(1/\varepsilon)}.
\]
Its area is comparable to $1/\log(1/\varepsilon)$ and tends to zero, whereas its diameter is comparable to one. At the origin its ratio to $h$ tends to infinity, so it has no uniform upper bound.
\end{proof}

Let $V\to C$ be a holomorphic vector bundle of positive rank with a fixed smooth Hermitian fiber norm. Write $U_\delta=\{v\in V:|v|<\delta\}$ and identify its zero section with $C$.

\begin{proposition}\label{prop:radial-tube-diameter}
Let $g_j$ be smooth Hermitian metrics on $U_{\delta_0}$. Suppose their restrictions to $C$ satisfy
\[
 g_j|_C\leq M h,\qquad
 A_j=\operatorname{Area}_{g_j|_C}(C)\longrightarrow0
\]
for a fixed $M>0$. Suppose also that there are $B>0$ and $0<\alpha\leq1$ such that, for every unit fiber vector $v$ over any $x\in C$, the ray $\gamma_{x,v}(s)=sv$ satisfies
\[
 g_j(\gamma_{x,v}'(s),\gamma_{x,v}'(s))
       \leq B s^{2\alpha-2},
       \qquad 0<s<\delta_0,
\]
uniformly in $j,x,v$. Then
\[
 \operatorname{diam}^{\mathrm{int}}(U_\delta,g_j)
 \leq\frac{2\sqrt B}{\alpha}\delta^\alpha
       +K_h\sqrt{A_j\bigl(1+\log(M/A_j)\bigr)},
       \qquad 0<\delta<\delta_0.
\]
In particular,
\[
 \lim_{\delta\downarrow0}\limsup_{j\to\infty}
       \operatorname{diam}^{\mathrm{int}}(U_\delta,g_j)=0.
\]
\end{proposition}
\begin{proof}
A point of fiber radius $r<\delta$ joins its base point on the zero section along the radial segment, entirely inside $U_\delta$. Its length is at most
\[
 \sqrt B\int_0^r s^{\alpha-1}\,\mathrm{d}s
       =\frac{\sqrt B}{\alpha}r^\alpha.
\]
The metric is smooth at the zero section for each $j$; the integrable uniform bound controls the whole segment. Join the two base points by an almost minimizing path in $(C,g_j|_C)$. This path also lies in $U_\delta$, and \cref{lem:conformal-area-diameter} bounds its length after its arbitrarily small error tends to zero. The two radial segments give the displayed intrinsic-diameter bound. Since $A_j(1+\log(M/A_j))\to0$, taking $j\to\infty$ and then $\delta\downarrow0$ proves the final assertion.
\end{proof}

For the node model, the exceptional area identity supplies $A_j\to0$, but a uniform upper metric bound on the exceptional curve and the radial estimate remain additional hypotheses. The proposition verifies the approximating-neighborhood condition of \cref{prop:finite-collapse-distance} when those two bounds hold; it does not verify the punctured limiting-neighborhood condition or conical asymptotics.
